\chapter{S07. Materiales III: tensiones, fluencia, difusión e inspección}
\noindent\textbf{Tiempo orientativo: 100 minutos. Cuatro preguntas; 10 puntos.}
\section*{1. Estado plano de tensiones (2.5 puntos)}
Un punto de una chapa ficticia soporta $\sigma_x=120\,\mathrm{MPa}$, $\sigma_y=-40\,\mathrm{MPa}$ y $\tau_{xy}=60\,\mathrm{MPa}$. Calcule las tensiones principales, la tensión tangencial máxima en el plano, la orientación del plano principal mayor y la tensión equivalente de von Mises. Con un límite elástico docente de $350\,\mathrm{MPa}$, indique el coeficiente de seguridad y qué no cubre.

\section*{2. Parámetro de Larson-Miller (2.5 puntos)}
Una pieza ficticia rompe a 100 h a 1000 K y se adopta $C=20$ con el tiempo en horas. Calcule el parámetro y el tiempo de rotura a 1050 K con la misma tensión. Explique dos límites del método.

\section*{3. Profundidad de cementación (2.5 puntos)}
El carbono difunde en austenita con $D_0=2.3\times10^{-5}\,\mathrm{m^2/s}$ y $Q=148\,\mathrm{kJ/mol}$. Calcule $D$ a 1223 K y la profundidad característica $\sqrt{Dt}$ tras 4 h, con $R=8.314462618\,\mathrm{J/(mol\,K)}$. Explique por qué no es un perfil de dureza.

\section*{4. Ultrasonidos por pulso-eco (2.5 puntos)}
Con $v=6300\,\mathrm{m/s}$ en aluminio, calcule el tiempo de ida y vuelta en una pared de 25 mm y la profundidad de un eco recibido a $4.76\,\mu\mathrm s$. Calcule la diferencia en dB entre dos ecos con razón de amplitudes 3 y explique por qué una amplitud no determina el tamaño de la discontinuidad.

\chapter{S08. Materiales IV: fatiga, daño acumulado, compuestos y fractura}
\noindent\textbf{Tiempo orientativo: 110 minutos. Cuatro preguntas; 10 puntos.}
\section*{1. Curva de Basquin (2.5 puntos)}
Un material ficticio cumple $\sigma_a=\sigma_f'(2N_f)^b$ con $\sigma_f'=900\,\mathrm{MPa}$ y $b=-0.09$. Calcule los ciclos hasta rotura para $\sigma_a=300\,\mathrm{MPa}$ y explique dos límites de la ecuación.

\section*{2. Regla de Miner (2.5 puntos)}
Un bloque de carga contiene $2\times10^4$ ciclos de una amplitud con vida $10^5$, $5\times10^4$ de otra con vida $5\times10^5$ y $10^5$ de una tercera con vida $2\times10^6$. Calcule el daño por bloque y los bloques hasta $D=1$. Explique por qué esa predicción no garantiza una vida segura.

\section*{3. Lámina unidireccional (2.5 puntos)}
Una lámina tiene $V_f=0.55$, $E_f=200\,\mathrm{GPa}$, $E_m=4\,\mathrm{GPa}$, $\rho_f=1800\,\mathrm{kg/m^3}$ y $\rho_m=1200\,\mathrm{kg/m^3}$. Calcule $E_L$, $E_T$ por el modelo de Reuss, la densidad y el módulo longitudinal específico. Explique por qué la diferencia entre $E_L$ y $E_T$ obliga a laminar en varias orientaciones.

\section*{4. Resistencia residual (2.5 puntos)}
Una placa con grieta superficial tiene $K_{IC}=40\,\mathrm{MPa\sqrt m}$ y $Y=1.12$. Calcule la tensión de fallo para semilongitudes de 4 mm y 8 mm, y el tamaño crítico a $250\,\mathrm{MPa}$. Explique la utilidad de la curva de resistencia residual.

\chapter{S09. Mantenimiento III: contadores, vida límite, fiabilidad y tendencia}
\noindent\textbf{Tiempo orientativo: 90 minutos. Cuatro preguntas; 10 puntos.}
Todos los datos son ficticios y las reglas de intervalo se dan en el enunciado; no corresponden a ningún programa aprobado.
\section*{1. Intervalo con dos contadores (2.5 puntos)}
Una tarea ficticia se repite cada 200 FH o 12 meses, lo que ocurra antes. Se realizó el 2 de marzo de 2026 a 3100 FH. El 8 de octubre de 2026 la aeronave tiene 3250 FH y vuela una media de 6 FH al día. Calcule el margen de cada contador y señale cuál gobierna.

\section*{2. Pieza de vida limitada (2.5 puntos)}
Una pieza tiene vida límite de 20000 ciclos y acumula 17450; se suman 5 ciclos al día. El plazo de aprovisionamiento es de 90 días. Calcule los ciclos y días restantes y cuántos días quedan hasta el último momento para pedir el repuesto.

\section*{3. Fiabilidad básica (2.5 puntos)}
Se registraron 6 retiradas no programadas de un componente en 18000 FH. Con tasa constante, calcule la tasa de retirada, el tiempo medio entre retiradas, la probabilidad de no retirada en 500 FH y las retiradas esperadas al mes en una flota de 12 aeronaves con 40 FH cada una. Explique la hipótesis de tasa constante.

\section*{4. Tendencia de un parámetro de motor (2.5 puntos)}
El margen de temperatura de gases de escape de un motor ficticio es 38 K y desciende linealmente $0.9\,\mathrm K$ cada 100 FH. Se define una alerta a 15 K. Calcule las horas y los días de uso (6 FH al día) hasta la alerta y explique dos limitaciones de la extrapolación lineal.

\chapter{S10. Combinada: fluencia, tensión térmica, corrosión y propagación}
\noindent\textbf{Tiempo orientativo: 110 minutos. Cuatro preguntas; 10 puntos.}
\section*{1. Fluencia de Norton (2.5 puntos)}
Una pieza ficticia tiene velocidad de fluencia estacionaria $10^{-8}\,\mathrm{s^{-1}}$ a 150 MPa y 1000 K, con exponente $n=5$. Calcule el cociente de velocidades al subir la tensión a 180 MPa, la velocidad nueva y el tiempo hasta 1 \% de deformación en ambos casos, a temperatura constante. Indique qué fase de fluencia no se representa.

\section*{2. Tensión térmica (2.5 puntos)}
Una barra con $E=200\,\mathrm{GPa}$ y $\alpha=13\times10^{-6}\,\mathrm{K^{-1}}$ sufre un calentamiento de 250 K con coacción total. Calcule la deformación libre, la tensión con coacción total y con coacción del 50 \%. Compare con un límite elástico docente de $400\,\mathrm{MPa}$.

\section*{3. Velocidad de corrosión (2.5 puntos)}
En un ensayo ficticio, una probeta de aluminio ($2.7\,\mathrm{g/cm^3}$, $25\,\mathrm{cm^2}$) pierde $0.84\,\mathrm g$ en 480 h. Calcule la velocidad de penetración en mm/año y explique por qué no puede extrapolarse sin más al servicio.

\section*{4. Propagación de grieta (2.5 puntos)}
Integre la ley de Paris con $C=3\times10^{-12}$ (m/ciclo, $\Delta K$ en $\mathrm{MPa\sqrt m}$), $m=3$, $Y=1.12$ y $\Delta\sigma=90\,\mathrm{MPa}$ entre semilongitudes de 2 mm y 9 mm. Indique dos hipótesis.
